% !TeX root = ../../Book3.tex
% 纲要标题：### 19.2 结构性质
% 纲要位置：计划纲要.md，第 2094 行。

\section{结构性质}
\label{b3:sec:1902}


% 纲要内容：
%
% * 商去正则元素
% * 局部化及适当的完备化性质
% * 正则局部环是 Cohen–Macaulay 环
%

\subsection{商去正则元素}
\label{b3:subsec:1902:01}

\begin{theorem}\label{b3:thm:cm-regular-quotient}
设 \((R,\mathfrak m)\) 为交换 Noether 局部环，\(M\ne0\) 为有限 \(R\) 模，\(x_1,\ldots,x_r\in\mathfrak m\) 为 \(M\) 正则序列。令 \(\overline M=M/(x_1,\ldots,x_r)M\)。则
\[
M\text{ 为 CM 模}\quad\Longleftrightarrow\quad
\overline M\text{ 为 CM 模},
\qquad \dim\overline M=\dim M-r.
\]
右边既可把 \(\overline M\) 看成 \(R\) 模，也可看成 \(R/(x_1,\ldots,x_r)\) 模。
\end{theorem}
\begin{proof}
各个逐次商非零，因为 \(M\) 有限且所有元素属于极大理想。对每一步，\cref{b3:cor:regular-quotient-depth}使深度恰减一，\cref{b3:prop:regular-quotient-dimension}使维数恰减一，故差 \(\dim M-\operatorname{depth}M\) 不变。商环和原环作用给出的子模、支集链及正则乘法完全相同，因而不会改变所比较的深度和维数。
\end{proof}

非零因子假设不能删除。取 \(R=k[x,y]_{(x,y)}/(x^2,xy)\)，则商去 \(x\) 后得到正则的一维局部环 \(k[y]_{(y)}\)，而原环的非零元 \(x\) 被整个极大理想零化，深度为零。这时商环的 CM 性不能反映原环的 CM 性。

\subsection{Cohen–Macaulay 性的局部化与完备化}
\label{b3:subsec:1902:02}

\begin{theorem}\label{b3:thm:cm-localization}
设 \((R,\mathfrak m)\) 为交换 Noether 局部环，\(M\ne0\) 为有限 CM 模。对每个 \(\mathfrak p\in\operatorname{Supp}_RM\)，\(M_{\mathfrak p}\) 都是 \(R_{\mathfrak p}\) 上的 CM 模。因而一般 Noether 环上有限模的 CM 性可只在极大理想处检查，并在任意局部化后保持。
\end{theorem}
\begin{proof}
在支集内，把 \(\mathfrak p\) 以下的一条最长链与 \(\mathfrak p\) 到 \(\mathfrak m\) 的一条最长链拼接，得
\(\dim M_{\mathfrak p}+\dim R/\mathfrak p\le\dim M\)。另一方面，\cref{b3:thm:depth-localization}给出
\[
\dim M=\operatorname{depth}_RM
\le\operatorname{depth}_{R_{\mathfrak p}}M_{\mathfrak p}+\dim R/\mathfrak p
\le\dim M_{\mathfrak p}+\dim R/\mathfrak p
\le\dim M.
\]
所以所有不等式都是等式，尤其局部化后的深度等于维数。一般情形中，给定素理想先把它包含在一个极大理想中，再应用已证的局部结论。局部化环的素点又与原环中避开乘法集的素点对应，故最后一项结论也成立。
\end{proof}

\begin{theorem}\label{b3:thm:cm-completion}
设 \((R,\mathfrak m)\) 为交换 Noether 局部环，\(M\ne0\) 为有限模。记 \(\widehat R\) 和 \(\widehat M\) 为极大理想进完备化，则
\[
\operatorname{depth}_{\widehat R}\widehat M=\operatorname{depth}_RM,
\qquad M\text{ 为 CM 模}\Longleftrightarrow\widehat M\text{ 为 CM 模}.
\]
\end{theorem}
\begin{proof}
令 \(k=R/\mathfrak m\)。取 \(k\) 的有限秩自由消解；它可以无限延伸，但每项有限，因为 \(R\) 为 Noether 环。\(\widehat R\) 的平坦性使张量后的复形仍消解同一个剩余域。有限秩保证逐项有
\[
\operatorname{Hom}_R(F_i,M)\otimes_R\widehat R
\cong\operatorname{Hom}_{\widehat R}(F_i\otimes_R\widehat R,\widehat M).
\]
平坦张量保持核、像及其商，因此上同调给出
\[
\operatorname{Ext}_R^i(k,M)\otimes_R\widehat R
\cong\operatorname{Ext}_{\widehat R}^i(k,\widehat M).
\]
忠实平坦性检测左侧是否为零。由\cref{b3:thm:depth-ext}，两边首次非零的次数相同，所以深度相同。\cref{b3:thm:completion-dimension}已证明有限模的维数也相同，遂得 CM 性的等价。
\end{proof}

这里使用的完备化是 Noether 局部环的极大理想进完备化。若在任意非局部环中选一个理想完备化，映射未必忠实；它可能完全丢掉某些素点，因而不能不加条件地用完备环检测原环的全部 CM 性。

\subsection{正则局部环的 Cohen–Macaulay 性}
\label{b3:subsec:1902:03}

\begin{corollary}\label{b3:cor:regular-ring-cm}
每个交换 Noether 正则局部环都是 CM 环。若交换 Noether 环的所有素点局部环均正则，则该环为 CM 环。
\end{corollary}
\begin{proof}
\cref{b3:thm:regular-local-characterizations}说明：一个 \(d\) 维正则局部环的极大理想可由长度为 \(d\) 的正则序列生成。因此深度至少为 \(d\)，结合深度不超过维数即得相等。非局部结论按定义逐点应用。
\end{proof}

\(k[[x_1,\ldots,x_d]]\) 及 \(k[x_1,\ldots,x_d]_{(x_1,\ldots,x_d)}\) 中，变量依次正则且最后的商为 \(k\)。把这些环商去一个非零非单位，就得到维数减一的 CM 环，但商环可能有奇点，甚至不约化。下一节将统一处理由多个正则元素给出的商。
